\chapter{Pruebas y comparativa}
\label{cap:comparativa}

Con el fin de contextualizar los resultados del capítulo~\ref{cap:explotacion},
este capítulo compara, por un lado, la metodología propia derivada en el
capítulo~\ref{cap:metodologia} frente a la OWASP WSTG completa de la que
parte; y, por otro, repite la explotación de vulnerabilidades equivalentes sobre dos
laboratorios de referencia ampliamente utilizados en formación en
ciberseguridad: DVWA (\emph{Damn Vulnerable Web Application}) y OWASP Juice
Shop, comparando los tres entornos desde varios criterios objetivos.

\section{Comparativa metodológica: proceso propio frente a la OWASP WSTG}
\label{sec:comparativa-metodologica}

La tabla~\ref{tab:comparativa-metodologica} compara la metodología propia
aplicada en el capítulo~\ref{cap:explotacion} con la OWASP WSTG completa de
la que se deriva (apartado~\ref{sec:derivacion-metodologia}), en términos de
cobertura y de esfuerzo estimado de aplicación.

\begin{table}[htbp]
  \centering
  \small
  \begin{tabular}{p{3.0cm}p{5.3cm}p{5.3cm}}
    \toprule
    \tabhead{Criterio} & \tabhead{OWASP WSTG completa} & \tabhead{Metodología propia} \\
    \midrule
    Categorías cubiertas &
      11 categorías, más de 90 pruebas individuales &
      6 pruebas, correspondientes a las categorías del OWASP Top~10 de mayor prevalencia e impacto \\
    Pasos por prueba &
      Variable según categoría, con sub-pasos específicos de cada prueba &
      Ciclo único y fijo de 6 pasos, aplicado igual a toda vulnerabilidad \\
    Perfil requerido &
      Equipo de auditoría con conocimiento amplio en todas las categorías &
      Equipo de desarrollo con conocimiento del OWASP Top~10 \\
    Esfuerzo estimado &
      Alto; propio de una auditoría externa de alcance general &
      Bajo-medio; aplicable de forma recurrente por un equipo interno \\
    Contexto de uso recomendado &
      Auditoría formal completa, con contrato y alcance definido &
      Validación ágil y recurrente en equipos pequeños o pymes con recursos limitados \\
    \bottomrule
  \end{tabular}
  \caption{Comparativa entre la OWASP WSTG completa y la metodología propia.}
  \label{tab:comparativa-metodologica}
\end{table}

La metodología propia cubre así una fracción reducida de la WSTG completa (6
de más de 90 pruebas), pero concentrada en las categorías de mayor
prevalencia según el OWASP Top~10~\citelib{OWASP21}, lo que permite obtener,
con una fracción del esfuerzo, una primera valoración de los riesgos más
probables de una aplicación. No sustituye a una auditoría completa, pero
resulta más ágil para su aplicación recurrente por un equipo con recursos
limitados, que es precisamente el objetivo con el que se derivó
(apartado~\ref{sec:derivacion-metodologia}). No se ha medido el tiempo real
de aplicación de la WSTG completa sobre \breachlab, por lo que la
comparación de esfuerzo anterior es una estimación cualitativa y no una
medición cronometrada; una comparación cuantitativa queda propuesta como
trabajo futuro (apartado~\ref{sec:trabajo-futuro}).

\section{Comparativa de entornos: DVWA y OWASP Juice Shop}

\subsection{Vulnerabilidades equivalentes en DVWA y Juice Shop}

DVWA reproduce, entre otras, inyección SQL, XSS almacenado y reflejado,
CSRF y control de acceso roto, con distintos niveles de dificultad
seleccionables (\texttt{low}, \texttt{medium}, \texttt{high},
\texttt{impossible}) que permiten observar cómo evoluciona una misma
contramedida. OWASP Juice Shop, por su parte, es una aplicación
Angular/Node.js vulnerable de forma intencionada que organiza sus retos
(«\emph{challenges}») en categorías alineadas con el OWASP Top~10, el OWASP
ASVS y el OWASP Automated Threat Handbook. La tabla~\ref{tab:equivalencias}
recoge, para cada vulnerabilidad de \breachlab, el módulo o reto
equivalente en cada laboratorio, como referencia para quien desee reproducir
la comparación de forma práctica.

\begin{table}[htbp]
  \centering
  \small
  \begin{tabular}{lll}
    \toprule
    \tabhead{Vulnerabilidad en \breachlab} & \tabhead{Módulo en DVWA} & \tabhead{Reto en Juice Shop} \\
    \midrule
    Inyección SQL             & SQL Injection / SQL Injection (Blind) & Login Admin \\
    XSS basado en DOM         & XSS (DOM)                             & DOM XSS \\
    CSRF                      & CSRF                                  & CSRF \\
    Fallos de autenticación   & Brute Force                           & Sin reto directamente equivalente \\
    IDOR                      & Sin módulo directamente equivalente   & View Basket / Manipulate Basket \\
    Configuración insegura    & Sin módulo directamente equivalente   & CSP Bypass / HTTP-Header XSS \\
    \bottomrule
  \end{tabular}
  \caption{Módulos y retos de DVWA y OWASP Juice Shop equivalentes a cada
    vulnerabilidad de \breachlab.}
  \label{tab:equivalencias}
\end{table}

\subsection{Comparativa}

\begin{table}[htbp]
  \centering
  \small
  \begin{tabular}{p{3.0cm}p{2.6cm}p{3.1cm}p{3.4cm}}
    \toprule
    \tabhead{Criterio} & \tabhead{DVWA} & \tabhead{OWASP Juice Shop} & \tabhead{BreachLab} \\
    \midrule
    Arquitectura &
      Monolito PHP + MySQL &
      SPA Angular + API Node.js &
      SPA Astro + API Flask (mismo origen) \\
    Cobertura de vulnerabilidades &
      Amplia, con niveles de dificultad &
      Muy amplia (todo el OWASP Top~10) &
      Acotada a 6 vulnerabilidades, en profundidad \\
    Rama segura incluida &
      Sí (nivel \texttt{impossible}) &
      No (hay que corregir aparte) &
      Sí, en el mismo código (\texttt{?safe=1}) \\
    Evidencia objetiva automatizada &
      No &
      Puntuación de retos &
      Registro de eventos (\texttt{attacks.log}) \\
    Valor didáctico para este TFM &
      Referencia de amplitud &
      Referencia de amplitud y de retos &
      Control total del código y de la comparación vulnerable/seguro \\
    \bottomrule
  \end{tabular}
  \caption{Comparativa entre DVWA, OWASP Juice Shop y BreachLab.}
  \label{tab:comparativa}
\end{table}

Como ampliación objetiva de esta comparativa, se recomienda ejecutar un
análisis DAST automatizado (por ejemplo, un escaneo \emph{baseline} de OWASP
ZAP~\citelink{ZAP}) contra los tres laboratorios, en su modo vulnerable y, cuando exista,
en su modo corregido, y recoger en una tabla el número de alertas por
severidad detectadas en cada caso.

\subsection{Conclusión de la comparativa}

DVWA y OWASP Juice Shop resultan más adecuados para adquirir una visión
amplia de vulnerabilidades diversas, pero al ser aplicaciones de
terceros no permiten instrumentar el código para generar evidencias propias
ni modificar la arquitectura para estudiar su efecto sobre una
vulnerabilidad concreta. \breachlab, al ser un desarrollo propio con
arquitectura SPA + API y ambas ramas de código en el mismo lugar, resulta
más limitado en cobertura pero permite un análisis más profundo y
verificable de cada vulnerabilidad.
